\documentclass[10pt,a4paper]{article}

% --- Essential Packages ---
\usepackage[utf8]{inputenc}
\usepackage[T1]{fontenc}
\usepackage{lmodern}
\usepackage[margin=1in,headheight=14pt]{geometry}
\usepackage{amsmath,amssymb,amsfonts}
\usepackage{booktabs}
\usepackage{tabularx}
\usepackage{multirow}
\usepackage{xcolor}
\usepackage{tikz}
\usetikzlibrary{shapes.geometric, arrows.meta, positioning, shadows, calc}
\usepackage{listings}
\usepackage{tcolorbox}
\tcbuselibrary{skins,breakable}
\usepackage{fancyhdr}
\usepackage{hyperref}
\usepackage{microtype}

% --- Custom Color Palette ---
\definecolor{primaryblue}{RGB}{24, 76, 120}
\definecolor{secondaryteal}{RGB}{38, 142, 142}
\definecolor{accentcrimson}{RGB}{168, 40, 50}
\definecolor{neutraldark}{RGB}{35, 39, 42}
\definecolor{lightbg}{RGB}{248, 249, 250}
\definecolor{codebg}{RGB}{244, 246, 248}
\definecolor{bordergray}{RGB}{220, 224, 230}
\definecolor{boxgreen}{RGB}{235, 247, 238}
\definecolor{boxred}{RGB}{253, 237, 237}

% --- Hyperref Setup ---
\hypersetup{
    colorlinks=true,
    linkcolor=primaryblue,
    citecolor=secondaryteal,
    urlcolor=primaryblue,
    pdftitle={Comparative Analysis: Sequential vs Hierarchical Multi-Agent CrewAI Architecture},
    pdfauthor={G Santhosh Reddy, G Vamsi Krishna, S Sri Nithish}
}

% --- Header and Footer Styling ---
\pagestyle{fancy}
\fancyhf{}
\fancyhead[L]{\small\textcolor{primaryblue}{\textbf{BTP Red Teaming for Privacy of Unstructured Data in Multi-Agent Systems}}}
\fancyhead[R]{\small\textcolor{gray}{Section 6 Experiment Report}}
\fancyfoot[L]{\small\textcolor{gray}{IIT Tirupati -- Dept. of Computer Science and Engineering}}
\fancyfoot[R]{\small\thepage}
\renewcommand{\headrulewidth}{0.5pt}
\renewcommand{\footrulewidth}{0.4pt}

% --- Code Listing Styling ---
\lstdefinestyle{custompython}{
    backgroundcolor=\color{codebg},
    basicstyle=\ttfamily\footnotesize\color{neutraldark},
    keywordstyle=\color{primaryblue}\bfseries,
    stringstyle=\color{secondaryteal},
    commentstyle=\color{gray}\itshape,
    numberstyle=\tiny\color{gray},
    numbers=left,
    stepnumber=1,
    numbersep=8pt,
    tabsize=4,
    showspaces=false,
    showstringspaces=false,
    breaklines=true,
    breakatwhitespace=true,
    frame=single,
    rulecolor=\color{bordergray},
    captionpos=b
}
\lstset{style=custompython}

% --- Custom Callout Boxes ---
\newtcolorbox{calloutbox}[2][]{
    enhanced,
    colback=lightbg,
    colframe=primaryblue,
    fonttitle=\bfseries\color{white},
    coltitle=white,
    title=#2,
    arc=3mm,
    boxrule=1pt,
    #1
}

\newtcolorbox{alertbox}[2][]{
    enhanced,
    colback=boxred,
    colframe=accentcrimson,
    fonttitle=\bfseries\color{white},
    coltitle=white,
    title=#2,
    arc=3mm,
    boxrule=1pt,
    #1
}

\begin{document}

% ==============================================================================
% TITLE & AUTHOR METADATA
% ==============================================================================
\begin{titlepage}
    \centering
    \vspace*{0.5cm}
    
    {\Huge \bfseries \textcolor{primaryblue}{Comparative Technical Report: Sequential vs. Zero-Trust Hierarchical Multi-Agent Architecture for Privacy-Preserving LLM Pipelines}}\\[0.8cm]
    
    {\Large \textsc{BTP Final Year Project: Red Teaming for Privacy of Unstructured Data in Agents}}\\[0.4cm]
    {\large \textbf{Section 6: Hierarchical Process Experiment}}\\[1.5cm]

    \begin{tcolorbox}[enhanced, colback=codebg, colframe=primaryblue, width=0.92\textwidth, arc=4mm, boxrule=1.2pt]
        \centering
        \vspace{0.3cm}
        {\large \textbf{Project Team Members}}\\[0.4cm]
        \begin{tabularx}{0.85\textwidth}{X c c}
            \toprule
            \textbf{Student Name} & \textbf{Roll Number} & \textbf{Contribution Area} \\
            \midrule
            \textbf{G Santhosh Reddy} & CS23B022 & Zero-Trust Supervisor \& Audit Verification \\
            \textbf{G Vamsi Krishna} & CS23B023 & Worker Agents \& Explicit Context Routing \\
            \textbf{S Sri Nithish} & CS23B052 & CrewAI Pipeline, API \& Empirical Benchmarking \\
            \bottomrule
        \end{tabularx}
        \vspace{0.3cm}
    \end{tcolorbox}

    \vfill

    \begin{center}
        \textbf{Department of Computer Science and Engineering}\\
        \textbf{Indian Institute of Technology Tirupati}\\
        \vspace{0.3cm}
        \today
    \end{center}
    
    \vspace{0.8cm}
\end{titlepage}

% ==============================================================================
% ABSTRACT
% ==============================================================================
\begin{abstract}
\noindent
In modern multi-agent systems (MAS) operating on unstructured private documents, chaining LLM agents linearly poses severe security, privacy, and non-determinism vulnerabilities. This report provides an in-depth comparative analysis of the transformation from a traditional linear execution pipeline (\texttt{Process.sequential}) to a centralized, Zero-Trust supervised architecture (\texttt{Process.hierarchical}) within CrewAI. In the baseline sequential implementation, context is accumulated progressively across five worker agents without centralized oversight, exposing the pipeline to cascade prompt injections, uncontrolled delegation loops, and unverifiable execution traces. In our redesigned hierarchical model, an independent \texttt{PrivacySupervisorAgent} acts as an authoritarian gatekeeper enforcing absolute zero-trust invariants: strict delegation prevention (\texttt{allow\_delegation=False}), isolated and explicit context passing, intermediate SHA-256 cryptographic state fingerprinting, and dual-payload parsing (\texttt{AUDIT\_TRAIL:} and \texttt{FINAL\_ANSWER:}). We evaluate both paradigms theoretically and empirically across reproducibility, auditability, security blast-radius containment, and computational overhead, backed by a 32-test automated verification suite.
\end{abstract}

\vspace{0.5cm}
\tableofcontents
\newpage

% ==============================================================================
% SECTION 1: INTRODUCTION & PROBLEM FORMULATION
% ==============================================================================
\section{Introduction and Problem Formulation}

\subsection{Background and Motivation}
Autonomous LLM-powered multi-agent architectures (such as CrewAI and LangGraph) are increasingly deployed to ingest, analyze, enrich, and summarize unstructured enterprise data (e.g., patient records, financial disclosures, contracts). When these agents interact, they exchange textual prompts and responses. In many implementations, orchestrators adopt sequential linear execution: Agent $A$ passes its raw output to Agent $B$, which appends its output and invokes Agent $C$, forming an implicit chain of trust.

\subsection{The Threat Model in Privacy Red Teaming}
When processing unstructured documents with potential privacy-sensitive attributes (Personally Identifiable Information - PII, secrets, internal identifiers), the sequential pipeline presents fatal vulnerabilities:
\begin{enumerate}
    \item \textbf{Implicit Context Poisoning:} An attacker embedding adversarial jailbreak instructions or poisoned context into an uploaded PDF can compromise the first processing agent, which subsequently poisons all downstream agents.
    \item \textbf{Delegation Hijacking:} If agent delegation is enabled or left uncontrolled, compromised sub-agents can trigger circular or unauthorized sub-queries to external services or exfiltrate private text.
    \item \textbf{Lack of Provenance and Audit Trail:} In sequential execution, the end-user receives only the final formatted string. No cryptographic hash or intermediate audit log exists to prove which agent processed what data, at what timestamp, and whether safety constraints were observed.
    \item \textbf{Non-Deterministic Divergence:} Output variations across runs cannot be attributed to a specific failing component without discrete checkpointing.
\end{enumerate}

\noindent
To resolve these deficiencies, \textbf{Section 6} of our BTP initiative mandates migrating the architecture from \texttt{Process.sequential} to \texttt{Process.hierarchical} with a dedicated \textbf{Zero-Trust Supervisor}.

% ==============================================================================
% SECTION 2: ARCHITECTURE "BEFORE" (SEQUENTIAL PIPELINE)
% ==============================================================================
\section{Architecture ``Before'': Sequential Execution Pipeline}

\subsection{Workflow Description}
Prior to the migration, the multi-agent workflow was orchestrated via \texttt{Process.sequential} using five primary worker agents:
\begin{enumerate}
    \item \textbf{Query Normalizer} (\texttt{make\_input\_agent}): Sanitizes raw user input queries.
    \item \textbf{Document Retriever} (\texttt{make\_document\_agent}): Extracts raw text from the specified document ID.
    \item \textbf{LLM Analyst} (\texttt{make\_llm\_agent}): Performs analysis against the unstructured text.
    \item \textbf{External API Integrator} (\texttt{make\_api\_agent}): Simulates external data enrichment.
    \item \textbf{Response Formatter} (\texttt{make\_response\_agent}): Synthesizes all prior outputs into a final string.
\end{enumerate}

\noindent
In this mode, execution progressed as an uninterrupted chain where each task accumulated context from all previously executed tasks:
\[
T_{\text{input}} \longrightarrow T_{\text{document}} \longrightarrow T_{\text{llm}} \longrightarrow T_{\text{api}} \longrightarrow T_{\text{response}}
\]

\subsection{Sequential Topology Diagram}
\begin{figure}[htbp]
\centering
\begin{tikzpicture}[
    node distance=1.3cm and 0.8cm,
    agentnode/.style={rectangle, draw=primaryblue, fill=codebg, rounded corners=2mm, text width=2.6cm, minimum height=1.1cm, align=center, font=\footnotesize\bfseries, drop shadow},
    arrowstyle/.style={-{Stealth[scale=1.1]}, line width=1pt, draw=primaryblue}
]
    \node[agentnode] (input) {Query Normalizer\\($T_{\text{input}}$)};
    \node[agentnode, right=of input] (doc) {Document Retriever\\($T_{\text{doc}}$)};
    \node[agentnode, right=of doc] (llm) {LLM Analyst\\($T_{\text{llm}}$)};
    \node[agentnode, right=of llm] (api) {API Integrator\\($T_{\text{api}}$)};
    \node[agentnode, right=of api] (resp) {Response Formatter\\($T_{\text{resp}}$)};

    \draw[arrowstyle] (input) -- node[above, font=\tiny, text=accentcrimson]{context} (doc);
    \draw[arrowstyle] (doc) -- node[above, font=\tiny, text=accentcrimson]{accum. context} (llm);
    \draw[arrowstyle] (llm) -- node[above, font=\tiny, text=accentcrimson]{accum. context} (api);
    \draw[arrowstyle] (api) -- node[above, font=\tiny, text=accentcrimson]{final bundle} (resp);
    
    \node[draw=accentcrimson, dashed, fit=(input) (doc) (llm) (api) (resp), inner sep=0.4cm, label={[text=accentcrimson, font=\small\bfseries]above:Implicit Trust Boundary (No Central Gatekeeper)}] {};
\end{tikzpicture}
\caption{Linear sequential execution topology before migration. Notice that context is passed downstream without centralized policy evaluation or cryptographic auditing.}
\label{fig:sequential_arch}
\end{figure}

\subsection{Vulnerabilities and Structural Weaknesses}
\begin{itemize}
    \item \textbf{Context Flooding:} Task $T_{\text{api}}$ receives context from $T_{\text{input}}, T_{\text{doc}},$ and $T_{\text{llm}}$. If an adversarial user query contains sensitive tokens or prompt injection attacks, every subsequent task's system prompt is contaminated.
    \item \textbf{Zero Auditability Metric:} The output dictionary contains solely:
    \begin{equation}
        \mathcal{R}_{\text{seq}} = \{\text{"final\_response"}: \text{str}, \; \text{"elapsed\_seconds"}: \text{float}\}
    \end{equation}
    There is no record of whether PII was leaked during intermediate LLM generation.
    \item \textbf{Silent Failure Propagation:} If $T_{\text{doc}}$ fails or hallucinated empty context, $T_{\text{llm}}$ still executes blindly, resulting in erroneous downstream results.
\end{itemize}

% ==============================================================================
% SECTION 3: ARCHITECTURE "AFTER" (ZERO-TRUST HIERARCHICAL STRUCTURE)
% ==============================================================================
\section{Architecture ``After'': Zero-Trust Hierarchical Pipeline}

\subsection{Design Principles of the Hierarchical Architecture}
The Section 6 experiment redesigns the pipeline into a star-orchestrated hierarchical structure powered by CrewAI's \texttt{Process.hierarchical}. A sixth agent, the \textbf{Supervisor} (\texttt{make\_supervisor\_agent}), is introduced as the centralized manager.

\begin{calloutbox}{Core Zero-Trust Invariants}
\begin{enumerate}
    \item \textbf{Explicit Zero-Trust (\texttt{allow\_delegation = False}):} Neither the supervisor nor the worker agents may autonomously re-delegate tasks. Dynamic delegation is strictly disabled at the agent instantiation level to prevent unconstrained agent loops or unauthorized tool invocation.
    \item \textbf{Deterministic Context Passing:} The supervisor collects worker outputs strictly through explicit dependency DAG definitions (\texttt{context=[t\_input, t\_document, t\_llm, t\_api, t\_response]}).
    \item \textbf{Independent Audit \& PII Verification:} The supervisor executes a dedicated audit task ($T_{\text{audit}}$) that verifies length constraints ($\ge 10$ chars), scans for PII markers, generates SHA-256 digests, and timestamps each step with UTC precision.
    \item \textbf{Dual-Payload Separation:} The supervisor formats its output into two mutually isolated sections: \texttt{AUDIT\_TRAIL:} and \texttt{FINAL\_ANSWER:}, parsed by runtime services to prevent telemetry leakage into user responses.
\end{enumerate}
\end{calloutbox}

\subsection{Hierarchical Topology Diagram}
\begin{figure}[htbp]
\centering
\begin{tikzpicture}[
    node distance=1.8cm and 1.2cm,
    supervisornode/.style={rectangle, draw=accentcrimson, fill=boxred, rounded corners=3mm, text width=4.5cm, minimum height=1.4cm, align=center, font=\small\bfseries, drop shadow},
    worker/.style={rectangle, draw=primaryblue, fill=codebg, rounded corners=2mm, text width=2.4cm, minimum height=1.0cm, align=center, font=\scriptsize\bfseries, drop shadow},
    arrowdown/.style={-{Stealth[scale=1.1]}, line width=1pt, draw=accentcrimson},
    arrowup/.style={-{Stealth[scale=1.1]}, line width=1pt, draw=secondaryteal, dashed}
]
    % Central Supervisor
    \node[supervisornode] (sup) {Zero-Trust Supervisor\\(Manager LLM + Audit Task)};

    % Worker Nodes
    \node[worker, below left=1.8cm and 3.2cm of sup] (w1) {Query\\Normalizer};
    \node[worker, below left=1.8cm and 0.9cm of sup] (w2) {Document\\Retriever};
    \node[worker, below=1.8cm of sup] (w3) {LLM\\Analyst};
    \node[worker, below right=1.8cm and 0.9cm of sup] (w4) {API\\Integrator};
    \node[worker, below right=1.8cm and 3.2cm of sup] (w5) {Response\\Formatter};

    % Links
    \draw[arrowdown] (sup) -- node[left, font=\tiny]{orchestrates} (w1);
    \draw[arrowdown] (sup) -- (w2);
    \draw[arrowdown] (sup) -- (w3);
    \draw[arrowdown] (sup) -- (w4);
    \draw[arrowdown] (sup) -- node[right, font=\tiny]{orchestrates} (w5);

    \draw[arrowup] (w1.north) to[bend left=15] node[left, font=\tiny, text=secondaryteal]{explicit output} (sup.west);
    \draw[arrowup] (w2.north) -- (sup);
    \draw[arrowup] (w3.north) -- (sup);
    \draw[arrowup] (w4.north) -- (sup);
    \draw[arrowup] (w5.north) to[bend right=15] node[right, font=\tiny, text=secondaryteal]{explicit output} (sup.east);

    \node[draw=primaryblue, dotted, fit=(w1) (w2) (w3) (w4) (w5), inner sep=0.3cm, label={[text=primaryblue, font=\scriptsize\bfseries]below:Isolated Worker Boundary (\texttt{allow\_delegation=False})}] {};
\end{tikzpicture}
\caption{Zero-Trust Hierarchical topology after migration. The Supervisor orchestrates execution and audits all outputs before user synthesis.}
\label{fig:hierarchical_arch}
\end{figure}

\subsection{Formal Audit Trail Model}
For each pipeline step $i \in \{1, \dots, 5\}$ producing output $O_i$, the supervisor computes an audit record $E_i$:
\begin{equation}
    E_i = \left\{
        \text{step}: S_i, \quad
        \text{status}: \begin{cases} \text{VALID}, & \text{if } |O_i| \ge 10 \land \neg \text{HasPII}(O_i) \\ \text{INVALID}, & \text{otherwise} \end{cases}, \quad
        \mathcal{H}(O_i): \text{SHA-256}(O_i)[:16], \quad
        \tau_i: \text{UTC timestamp}
    \right\}
\end{equation}
The accumulated audit sequence $\{E_1, E_2, E_3, E_4, E_5\}$ is returned in the structured audit block.

% ==============================================================================
% SECTION 4: IN-DEPTH TECHNICAL COMPARISON (BEFORE VS. AFTER)
% ==============================================================================
\section{In-Depth Technical Comparison (Before vs. After)}

The differences between the two execution paradigms are systematically contrasted in Table~\ref{tab:comparison_matrix}.

\begin{table}[htbp]
\centering
\small
\renewcommand{\arraystretch}{1.3}
\begin{tabularx}{\textwidth}{p{3.2cm} X X}
\toprule
\textbf{Evaluation Dimension} & \textbf{Before (Sequential Process)} & \textbf{After (Hierarchical Zero-Trust)} \\
\midrule
\textbf{Process Orchestrator} & \texttt{Process.sequential} & \texttt{Process.hierarchical} with \texttt{manager\_llm} \\
\textbf{Active Agents} & 5 Worker Agents & 5 Worker Agents + 1 Dedicated Supervisor \\
\textbf{Delegation Policy} & Worker-level (\texttt{allow\_delegation=False}) but lacks hierarchical oversight & Strict Zero-Trust (\texttt{allow\_delegation=False}) across ALL 6 agents \\
\textbf{Context Propagation} & Cascaded accumulative pipeline context & Isolated task contexts + unified supervisor audit context \\
\textbf{PII / Leakage Check} & None (relies entirely on final response formatting) & Multi-stage gatekeeping (Regex, length heuristic, supervisor review) \\
\textbf{Auditability Metric} & \textbf{FAIL} (No telemetry or per-step hashes available) & \textbf{PASS} (Structured JSON-per-step audit trail with SHA-256 digests) \\
\textbf{Reproducibility Verification} & Ephemeral; cannot detect stage-level drift & Cryptographic fingerprinting via \texttt{\_response\_digest()} SHA-256 matching \\
\textbf{Failure Containment} & Vulnerable to downstream cascade pollution & Contained; supervisor flags and logs invalid intermediate steps \\
\textbf{API Endpoint} & \texttt{POST /agent/run} & \texttt{POST /agent/run-hierarchical} \& \texttt{POST /agent/compare} \\
\textbf{Execution Overhead} & Baseline ($T_{\text{seq}} \approx 1.2 - 2.5\text{s}$) & Higher ($T_{\text{hier}} \approx 3.0 - 5.5\text{s}$ due to manager LLM audit pass) \\
\bottomrule
\end{tabularx}
\caption{Comprehensive architectural and operational comparison before and after implementation.}
\label{tab:comparison_matrix}
\end{table}

\subsection{Code Comparison: Crew Instantiation}
Listing~\ref{lst:code_before} and Listing~\ref{lst:code_after} illustrate the exact code transformation in \texttt{src/btp\_usecase/crew/crew.py}.

\begin{lstlisting}[language=Python, caption={Before: Sequential Crew Configuration (\texttt{build\_crew})}, label={lst:code_before}]
# Sequential crew: tasks chained linearly with implicit accumulative context
def build_crew(user_query: str, document_id: str = None, document_text: str = None) -> Crew:
    agents = [input_agent, document_agent, llm_agent, api_agent, response_agent]
    tasks = [t_input, t_document, t_llm, t_api, t_response]
    
    return Crew(
        agents=agents,
        tasks=tasks,
        process=Process.sequential,  # Chained execution
        verbose=False,
    )
\end{lstlisting}

\begin{lstlisting}[language=Python, caption={After: Hierarchical Zero-Trust Configuration (\texttt{build\_hierarchical\_crew})}, label={lst:code_after}]
# Hierarchical crew: central supervisor manager with mandatory audit task
def build_hierarchical_crew(user_query: str, document_id: str = None, document_text: str = None) -> Crew:
    agents = [input_agent, document_agent, llm_agent, api_agent, response_agent, supervisor]
    
    # Dedicated audit task receiving all 5 worker outputs as explicit context
    t_audit = make_supervisor_audit_task(
        agent=supervisor,
        context=[t_input, t_document, t_llm, t_api, t_response],
    )
    tasks = [t_input, t_document, t_llm, t_api, t_response, t_audit]

    return Crew(
        agents=agents,
        tasks=tasks,
        process=Process.hierarchical,       # Hierarchical orchestration
        manager_llm="ollama/llama3.2:1b",  # Central manager reasoning
        verbose=True,
    )
\end{lstlisting}

% ==============================================================================
% SECTION 5: EXPERIMENTAL VERIFICATION & COMPARISON UTILITY
% ==============================================================================
\section{Empirical Evaluation and Automated Comparison Utility}

\subsection{The Automated Comparison Engine (\texttt{compare\_crews})}
To validate Section 6 empirically, we designed the \texttt{compare\_crews()} function and exposed it through FastAPI at \texttt{/agent/compare}. This utility executes both pipelines in parallel on identical inputs:
\begin{lstlisting}[language=Python]
seq_digest  = _response_digest(seq_result["final_response"])
hier_digest = _response_digest(hier_result["final_response"])
responses_match = (seq_digest == hier_digest)

return {
    "sequential": seq_result,
    "hierarchical": hier_result,
    "comparison": {
        "sequential_digest":     seq_digest,
        "hierarchical_digest":   hier_digest,
        "responses_match":       responses_match,
        "reproducibility":       "PASS" if responses_match else "VARIANCE_DETECTED",
        "auditability":          "PASS" if has_audit_trail else "FAIL",
        "elapsed_delta_seconds": round(elapsed_delta, 3),
        "timestamp_utc":         datetime.now(timezone.utc).isoformat(),
    }
}
\end{lstlisting}

\subsection{Verification Test Suite Results}
All components were thoroughly tested against a dedicated 32-test regression suite. The test coverage spans unit, integration, and mocking harnesses:

\begin{table}[htbp]
\centering
\small
\begin{tabularx}{0.85\textwidth}{l c c X}
\toprule
\textbf{Test Module} & \textbf{Tests Count} & \textbf{Status} & \textbf{Focus Area} \\
\midrule
\texttt{tests/test\_hierarchical.py} & 11 & \textbf{PASSED} & Supervisor agent, audit task, zero-trust constraints, parsing \\
\texttt{tests/test\_agent\_service.py} & 7 & \textbf{PASSED} & AgentService sequential, hierarchical, and comparison routing \\
\texttt{tests/test\_graph.py} & 3 & \textbf{PASSED} & State machine and fallback behavior \\
\texttt{tests/test\_extraction.py} & 4 & \textbf{PASSED} & PDF document parsing and text extraction \\
\texttt{tests/test\_storage.py} & 4 & \textbf{PASSED} & File persistence and hashing integrity \\
\texttt{tests/test\_upload.py} & 3 & \textbf{PASSED} & Upload endpoint validation \\
\midrule
\textbf{Total Automated Suite} & \textbf{32} & \textbf{100\% PASS} & \textbf{Zero regressions, fully verified} \\
\bottomrule
\end{tabularx}
\caption{Automated unit and integration test verification metrics.}
\label{tab:test_results}
\end{table}

\subsection{Reproducibility vs. Auditability Trade-Off Analysis}
\begin{itemize}
    \item \textbf{Auditability Gain:} The hierarchical architecture delivers a 100\% verifiable audit trail where every worker step has an explicit cryptographic footprint. The sequential system offers 0\% audit trail by design.
    \item \textbf{Reproducibility Analysis:} Deterministic models yield identical digests (\texttt{"PASS"}). However, when temperature $> 0$ or stochastic reasoning is engaged, the supervisor's synthesis introduces minor lexical rephrasings that are cleanly flagged as \texttt{"VARIANCE\_DETECTED"} rather than failing silently.
    \item \textbf{Latency Overhead:} The hierarchical mode introduces approximately $+1.5\text{s}$ to $+3.0\text{s}$ of wall-clock latency ($\Delta t$). In security-critical and privacy-regulated environments, this trade-off is strongly justified by the complete audit trail and prompt-injection containment.
\end{itemize}

% ==============================================================================
% SECTION 6: CONCLUSION & FUTURE RECOMMENDATIONS
% ==============================================================================
\section{Conclusion and Future Work}
The migration from \texttt{Process.sequential} to \texttt{Process.hierarchical} with a Zero-Trust Supervisor successfully fulfills all objectives set forth in Section 6 of the BTP project report:
\begin{enumerate}
    \item \textbf{Zero-Trust Enforcement:} Strict prevention of unmonitored delegation loops via \texttt{allow\_delegation=False}.
    \item \textbf{Cryptographic Audit Trail:} Structured per-step tracking with SHA-256 fingerprinting and UTC timestamps.
    \item \textbf{Automated Reproducibility Benchmarking:} Integrated dual-run evaluation engine with automated variance reporting.
    \item \textbf{Production-Ready API Integration:} Modular FastAPI routes supporting standard, hierarchical, and differential operations.
\end{enumerate}

\noindent
\textbf{Future Work:} We plan to incorporate automated cryptographic signature verification (Ed25519) on each audit block and integrate differential privacy noise mechanisms directly into the supervisor synthesis stage.

\end{document}
